\chapter{Fainting}

\section*{Definition and Overview}
Fainting, known medically as syncope, is a sudden, brief loss of consciousness caused by a temporary drop in the blood supply to the brain. It usually lasts a few seconds to a minute, and the person recovers quickly once they are lying down. Fainting is common and is usually harmless, most often caused by a simple faint in response to emotion, pain, heat, or standing for long periods. However, fainting can also be caused by serious heart conditions, which makes assessment important, particularly in older people and when fainting happens during exercise.

\section*{Epidemiology}
Fainting is very common; most people faint at least once in their lives, and it accounts for a large number of emergency department visits. Simple faints are most common in young people, while fainting in older people is more likely to have a serious cause, such as heart rhythm problems or medication-related low blood pressure. Faints that occur during exercise or without warning are of particular concern.

\section*{Aetiology and Pathophysiology}
Fainting occurs when the brain is briefly starved of blood.
\begin{itemize}
    \item \textbf{Vasovagal syncope:} the simple faint, in which emotion, pain, fear, or standing triggers a sudden slowing of the heart and a drop in blood pressure
    \item \textbf{Situational syncope:} fainting during urination, coughing, or straining
    \item \textbf{Orthostatic hypotension:} a drop in blood pressure on standing, often from medicines, dehydration, or age
    \item \textbf{Heart conditions:} abnormal rhythms, valve disease, and heart attack can cause fainting
    \item \textbf{Other causes:} low blood sugar, seizures, and some medicines
    \item \textbf{The mechanism:} blood pressure falls, blood flow to the brain drops, and consciousness is lost until the blood supply is restored, usually by lying down
\end{itemize}

\section*{Clinical Features}
The features help distinguish a simple faint from a serious cause.
\begin{itemize}
    \item Light-headedness, nausea, sweating, and blurred vision before the faint
    \item A brief loss of consciousness, usually seconds to a minute
    \item A quick recovery once lying down
    \item Precipitants: heat, emotion, pain, standing, or fasting
    \item In a simple faint: warning symptoms, a slow recovery of colour, and no injury
    \item Concerning features: fainting during exercise, without warning, with palpitations, or with chest pain
    \item In older people: fainting may be due to medicines or heart problems
\end{itemize}

\section*{Differential Diagnosis}
Other conditions cause loss of consciousness or collapse.
\begin{itemize}
    \item Seizures, which involve jerking, tongue biting, and a slow recovery
    \item Heart rhythm problems and heart attack
    \item Low blood sugar in people with diabetes
    \item Stroke and other neurological events
    \item Falls without loss of consciousness
    \item Drug and alcohol effects
\end{itemize}

\section*{When to Seek Medical Care}
A first faint in an older person, a faint during exercise, a faint with palpitations or chest pain, or a faint causing injury should be assessed by a doctor. Recurrent faints also warrant assessment.

\textbf{Call 000 (or local emergency services) for:}
\begin{itemize}
    \item A faint during exercise
    \item Fainting with chest pain, palpitations, or breathlessness
    \item Fainting with no warning signs, especially in an older person
    \item Fainting followed by confusion, weakness, or difficulty speaking
    \item A faint in pregnancy or with severe bleeding
\end{itemize}

\section*{Diagnosis and Investigations}
The cause is found by a careful history and targeted tests.
\begin{itemize}
    \item \textbf{History:} the circumstances, warning signs, recovery, and any medicines
    \item \textbf{Examination:} blood pressure lying and standing, and heart examination
    \item \textbf{ECG:} to check the heart rhythm
    \item \textbf{Blood tests:} for anaemia, blood sugar, and other causes
    \item \textbf{Holter monitoring:} for heart rhythm over days to weeks
    \item \textbf{Other tests:} tilt testing, echocardiography, or specialist referral as needed
\end{itemize}

\section*{Management}
Management depends on the cause.
\begin{itemize}
    \item \textbf{Simple faints:} reassurance and education about warning signs and prevention
    \item \textbf{Orthostatic hypotension:} adjusting medicines, drinking fluids, and rising slowly
    \item \textbf{Heart conditions:} treatment of the rhythm or valve problem
    \item \textbf{Counter-pressure manoeuvres:} tensing the muscles at the first warning signs
    \item \textbf{Lifestyle:} staying hydrated, avoiding triggers, and not skipping meals
    \item \textbf{Driving advice:} people with recurrent or unexplained faints may need advice about driving
\end{itemize}

\section*{Complications}
\begin{itemize}
    \item Injuries from falls during the faint
    \item Fractures, particularly in older people
    \item Head injury
    \item The underlying heart condition being missed
    \item Anxiety about fainting and avoidance of activities
\end{itemize}

\section*{Prognosis and Outcomes}
The outlook for fainting is excellent when it is a simple faint, and most people can prevent attacks by recognising the warning signs and lying down. Fainting from heart conditions is more serious but is treatable once identified. The key is assessment to distinguish the harmless from the dangerous.

\section*{Prevention and Self-Care}
\begin{itemize}
    \item At the first warning signs, sit or lie down and raise the legs
    \item Stand up slowly, especially after rest
    \item Drink plenty of fluids and avoid skipping meals
    \item Avoid prolonged standing in heat
    \item Recognise and manage your triggers
    \item Seek medical advice for recurrent faints or concerning features
    \item Follow advice about driving if relevant
\end{itemize}

\section*{Sources and Further Reading}
The following reputable sources provide further information for healthcare students and patients seeking reliable, current advice about fainting.
\begin{itemize}
    \item Healthdirect Australia. \textit{Fainting.} https://www.healthdirect.gov.au/fainting
    \item Heart Foundation Australia. \textit{Syncope and fainting.} https://www.heartfoundation.org.au/
    \item St John Ambulance Australia. \textit{First aid for fainting.} https://stjohn.org.au/
    \item NHS. \textit{Fainting.} https://www.nhs.uk/conditions/fainting/
    \item Mayo Clinic. \textit{Syncope (fainting).} https://www.mayoclinic.org/diseases-conditions/syncope/
    \item MSD Manual. \textit{Syncope.} https://www.msdmanuals.com/
\end{itemize}
